\section{Problem\'atica}

Las peque\~nas y medianas empresas (PyMEs) representan un pilar fundamental de las econom\'ias
latinoamericanas; sin embargo, su acceso a herramientas de inteligencia de negocios y an\'alitica
de datos sigue siendo limitado. Seg\'un \citet{tovar2011bi}, los principales obst\'aculos para la
adopci\'on de soluciones de \textit{Business Intelligence} en este segmento empresarial son el
costo elevado de las plataformas, la falta de personal cualificado y la complejidad de las
integraciones con sistemas existentes.

Esta brecha tecnol\'ogica se profundiza cuando se considera que, a diferencia de las grandes
corporaciones, las PyMEs carecen de departamentos de TI dedicados capaces de gestionar
infraestructuras de datos complejas \citep{cordero2020sistemas}. La \citeauthor{oecd2023barreras}
(\citeyear{oecd2023barreras}) se\~nala adem\'as que la incompatibilidad entre sistemas heredados y
las nuevas plataformas tecnol\'ogicas constituye una barrera estructural que impide la
modernizaci\'on del tejido empresarial en econom\'ias emergentes.

Por otro lado, la complejidad percibida en la adopci\'on de tecnolog\'ias de la informaci\'on
constituye en s\'i misma un factor disuasorio \citep{nurqamarani2021complejidad}. Las interfaces
t\'ecnicas de las herramientas de an\'alitica requieren conocimientos especializados en SQL,
estad\'istica y ciencia de datos que no est\'an al alcance del perfil medio de los trabajadores
de una PyME. Esto se traduce en una subutilizaci\'on de los datos organizacionales y en la
incapacidad de extraer valor de la informaci\'on que ya poseen estas empresas.

\citet{wamba2017desafios} constatan que, incluso cuando las organizaciones logran implementar
soluciones anal\'iticas, la falta de una cultura basada en datos y la escasez de competencias
internas generan un retorno de inversi\'on insatisfactorio. En consecuencia, existe una necesidad
urgente de soluciones de bajo c\'odigo o lenguaje natural que permitan democratizar el acceso a
la informaci\'on estructurada, reduciendo la dependencia de intermediarios t\'ecnicos.